\documentclass[11pt]{article}

\usepackage[T1]{fontenc}
\usepackage{amsmath,amssymb,amsthm,mathtools}
\usepackage[margin=1in]{geometry}
\usepackage{microtype}
\usepackage[hidelinks]{hyperref}

\newtheorem{theorem}{Theorem}[section]
\newtheorem{proposition}[theorem]{Proposition}
\newtheorem{corollary}[theorem]{Corollary}
\newtheorem{lemma}[theorem]{Lemma}
\theoremstyle{remark}
\newtheorem{remark}[theorem]{Remark}

\newcommand{\F}{\mathbb F}
\newcommand{\C}{\mathbb C}
\newcommand{\E}{\mathcal E}
\newcommand{\1}{\mathbf 1}
\newcommand{\supp}{\operatorname{supp}}
\newcommand{\disttr}{D_{\mathrm{tr}}}

\hypersetup{
  pdftitle={A Robust Inverse Uncertainty Principle for the Boolean Cube,
    Formalized in Lean},
  pdfauthor={Marios Georgiou}
}

\title{A Robust Inverse Uncertainty Principle for the Boolean Cube,\\
Formalized in Lean}
\author{Marios Georgiou\\Galois, Inc.}
\date{August 2026}

\begin{document}
\maketitle

\begin{abstract}
An uncertainty principle says that a state cannot be sharply concentrated in
two incompatible descriptions at once. On $n$ bits, the two descriptions
considered here are the ordinary basis and the Walsh--Hadamard basis. The
exact support uncertainty inequality says that the numbers of nonzero
coordinates in these two bases have product at least $2^n$. Its equality
cases are uniform superpositions over shifted linear subspaces, with a
possible sign pattern. Miller recently gave a robust coefficient inequality
for this phenomenon over arbitrary finite abelian groups
\cite{Miller2026}. We prove a quantitative refinement for the Boolean cube
that allows small probability leakage outside prescribed sets $S$ and $T$.
If the leakage is at most $\varepsilon$ and
$|S||T|\leq(1+\delta)2^n$, then, in a local parameter range, both sets are
close to complementary shifted subspaces and the state is close to the
corresponding coset state. The dependence is explicit:
\[
 q=\max\{(1024\varepsilon)^{1/4},\,4\sqrt{\delta}\},
\]
with relative set errors at most $5q$ and $12q$, and squared Euclidean state
error at most $264q$. The constants are not optimized. The complete proof is
formalized in Lean 4 against Mathlib, with a small Mathlib-only public
specification separated from the implementation.
\end{abstract}

\section{Introduction}

The word ``uncertainty'' has several mathematical meanings. The common idea
is that a vector can be simple in one coordinate system or in an
incompatible coordinate system, but generally not in both. For an $n$-qubit
state, the two coordinate systems in this note are:
\begin{itemize}
  \item the \emph{computational basis}, labelled by $n$-bit strings; and
  \item the \emph{Hadamard basis}, obtained by applying a Hadamard gate to
  every qubit.
\end{itemize}
The change of coordinates between them is the Walsh Fourier transform.

The \emph{support} of a function is simply the set of inputs where it is
nonzero. A finite version of the Donoho--Stark support uncertainty principle
says that every nonzero function $f$ on a finite abelian group $G$ satisfies
\[
  |\supp f|\,|\supp\widehat f|\geq |G|.
\]
The hat means ``take the Fourier transform.'' Equality has a rigid form:
up to harmless symmetries, $f$ is constant on a shifted subgroup and zero
elsewhere. On bit strings, subgroups are linear subspaces and shifted
subgroups are affine subspaces, also called cosets.

The exact support is unstable under noise. An arbitrarily small perturbation
can give a vector full support in both bases. For quantum applications it is
more natural to specify sets $S$ and $T$ that capture almost all measurement
probability. The question addressed here is:

\begin{quote}
If $S$ and $T$ nearly attain the support-product lower bound and a state is
almost concentrated on them in complementary bases, must the sets and state
be close to a hidden-coset configuration?
\end{quote}

For the Boolean cube the answer is yes, quantitatively. The result below
separates the two sources of error: $\varepsilon$ measures probability
leakage and $\delta$ measures excess in the support product. This separation
is useful both mathematically and for reading the formal statement.

The qualitative phenomenon is not new. Exact uncertainty and its equality
cases have an extensive literature
\cite{DonohoStark1989,Meshulam2006,BonamiGhobber2013,FengHollmannXiang2019,
WigdersonWigderson2021}. A recent theorem of Miller controls Fourier overlap
through uniformity and linearity coefficients on arbitrary finite abelian
groups; its support corollary diagnoses near equality in the support-product
bound \cite{Miller2026}. Applying the general theorem to normalized
restrictions connects it directly to the concentration-set regime considered
here. The result in this note should therefore be read as a quantitative
refinement for $\F_2^n$, rather than as a separate general principle. It
trades group generality for symmetric-difference bounds for paired cosets,
complex state-vector control including phase, explicit constants, and a
machine-checked proof. Section~\ref{sec:literature} gives a precise
comparison.
Section~\ref{sec:continuous} compares this result with continuous Heisenberg
stability, where the corresponding rigid objects are Gaussian wave packets
rather than cosets.

\section{The result in plain language}\label{sec:plain}

An element of $\F_2^n$ is an $n$-bit string. Addition in $\F_2^n$ is bitwise
exclusive-or. A \emph{linear subspace} $H$ is a collection of bit strings
that contains the all-zero string and is closed under this addition. A
\emph{coset} $a+H$ is what one gets by adding a fixed string $a$ to every
element of $H$.

A normalized function $f:\F_2^n\to\C$ can be read as a quantum state:
$|f(x)|^2$ is the probability of observing $x$ in the computational basis.
After a Hadamard transform, the probability of observing $y$ is
$|\widehat f(y)|^2/2^n$. Thus the assumptions of the theorem say:
\begin{enumerate}
  \item computational-basis measurement lands in a set $S$ with probability
  at least $1-\varepsilon$;
  \item Hadamard-basis measurement lands in a set $T$ with probability at
  least $1-\varepsilon$; and
  \item the product $|S||T|$ is only a factor $1+\delta$ above the smallest
  value permitted by uncertainty.
\end{enumerate}

The conclusion reconstructs a hidden linear structure. It finds a subspace
$H$ such that $S$ is close to one shifted copy of $H$, while $T$ is close to
a shifted copy of the complementary space $H^\perp$. Here $H^\perp$ consists
of the bit strings having dot product zero with every element of $H$.
The theorem also says that the entire complex state $f$, not just the two
sets, is close to
\[
  \frac{1}{\sqrt{|H|}}
  \sum_{x\in a+H}(-1)^{b\cdot x}|x\rangle,
\]
up to one physically irrelevant global phase.

There are three numerical parameters because they play different roles:
\begin{itemize}
  \item $\varepsilon$ is lost probability mass;
  \item $\delta$ is excess in the set-size product; and
  \item $q$ is the accuracy promised by the conclusion.
\end{itemize}
The theorem is local: it assumes $0<q\leq1/100$ and applies when
\[
  \varepsilon\leq q^4/1024,\qquad \delta\leq q^2/16.
\]
Equivalently, one can compute $q$ from the two observed errors by taking
\[
  q=\max\{(1024\varepsilon)^{1/4},4\sqrt{\delta}\}.
\]
The fourth root means that the present quantitative bound becomes useful
only when the leakage is quite small. This is a limitation of the proved
estimate, not a claim that fourth-root behavior is optimal.

\section{Formal statement of the result}

\subsection{Normalization}

Fix
\[
  V=\F_2^n,\qquad N=|V|=2^n.
\]
For $f:V\to\C$ use the unnormalized Walsh transform
\[
  \widehat f(y)=\sum_{x\in V}(-1)^{x\cdot y}f(x).
\]
Thus Parseval's identity is
\[
  \sum_{y\in V}|\widehat f(y)|^2
    =N\sum_{x\in V}|f(x)|^2.
\]
For $A\subseteq V$, write
\[
  \E_A(f)=\sum_{x\in A}|f(x)|^2.
\]
The symmetric difference of two sets is denoted by $\triangle$.

For a subspace $H\leq V$, offsets $a,b\in V$, and a phase $c\in\C$ with
$|c|=1$, define the normalized modulated coset state
\begin{equation}\label{eq:coset-state}
  \phi_{H,a,b,c}(x)
  =\frac{c(-1)^{b\cdot x}}{\sqrt{|H|}}\,\1_{a+H}(x).
\end{equation}
Its Walsh transform is supported on $b+H^\perp$.

\subsection{Two-error theorem}

\begin{theorem}[Robust inverse uncertainty]\label{thm:main}
Let $S,T\subseteq V$ be nonempty, let $f:V\to\C$, and let
$\varepsilon,\delta,q\in\mathbb R$ satisfy
\[
  \varepsilon,\delta\geq0,\qquad
  0<q\leq\frac1{100},\qquad
  \varepsilon\leq\frac{q^4}{1024},\qquad
  \delta\leq\frac{q^2}{16}.
\]
Suppose
\begin{align}
  \E_V(f)&=1,\label{eq:normalization}\\
  \E_S(f)&\geq1-\varepsilon,\label{eq:primal-concentration}\\
  \E_T(\widehat f)&\geq(1-\varepsilon)N,
    \label{eq:dual-concentration}\\
  |S|\,|T|&\leq(1+\delta)N.\label{eq:product}
\end{align}
Then there are a subspace $H\leq V$, offsets $a,b\in V$, and a phase
$c\in\C$, $|c|=1$, such that
\begin{align}
  |S\triangle(a+H)|&\leq5q|S|,\label{eq:S-shape}\\
  |T\triangle(b+H^\perp)|&\leq12q|T|,\label{eq:T-shape}\\
  \sum_{x\in V}|f(x)-\phi_{H,a,b,c}(x)|^2
    &\leq264q.\label{eq:state-shape}
\end{align}
\end{theorem}

The role of $q$ is only to make the conclusion and the admissible error
region easy to read. It can be eliminated.

\begin{corollary}[Independent-error form]\label{cor:rho}
Under \eqref{eq:normalization}--\eqref{eq:product}, put
\[
  \rho(\varepsilon,\delta)
  =\max\{(1024\varepsilon)^{1/4},\,4\sqrt{\delta}\}.
\]
If $0<\rho(\varepsilon,\delta)\leq1/100$, then
\eqref{eq:S-shape}--\eqref{eq:state-shape} hold with
$q=\rho(\varepsilon,\delta)$.
\end{corollary}

\begin{proof}
The definition gives
$\varepsilon\leq\rho^4/1024$ and
$\delta\leq\rho^2/16$, so Theorem~\ref{thm:main} applies.
\end{proof}

\begin{remark}[Local and non-sharp constants]
The condition $\rho\leq1/100$ makes this a local stability theorem. The
constants $5$, $12$, and $264$, and the powers and numerical factors in
$\rho$, are outputs of a proof designed for transparency and formalization;
no optimality is claimed. The endpoint $\varepsilon=\delta=0$ is covered by
the classical exact equality theorem. It is separated from
Corollary~\ref{cor:rho} because that corollary uses a strictly positive
control scale.
\end{remark}

\section{Proof architecture}\label{sec:proof}

This section gives a conventional account of the proof. The formal
development cited in Section~\ref{sec:lean} contains all finite-sum,
rounding, and constant calculations. A reader interested mainly in the
statement or the quantum interpretation can skip this section.

\subsection{The proof in five steps}

The argument can be summarized without the constants as follows.
\begin{enumerate}
  \item Restrict $f$ to $S$. The small amount removed in this operation
  remains small after the Walsh transform by Parseval's identity, which says
  that the transform preserves total squared length up to the known factor
  $2^n$.
  \item Near equality in the support bound makes the submatrix of Walsh
  signs with rows in $T$ and columns in $S$ close to a rank-one matrix. A
  rank-one matrix is one whose rows are all scalar multiples of a single
  row.
  \item A four-corner identity for Walsh signs converts this approximate
  rank-one property into a statement that most translated pairs from
  $S\times T$ have dot product zero.
  \item An additive-combinatorial argument shows that two nearly maximally
  sized, nearly orthogonal sets must be close to complementary subspaces.
  \item Projecting onto the two resulting cosets leaves little error. The
  common intersection of the two projector ranges is a single coset-state
  line, so $f$ is close to a unit vector on that line.
\end{enumerate}

\subsection{The restricted inverse proposition}

Let $P_S f=f\1_S$. The structural core of the proof is the following
restricted-Walsh statement.

\begin{proposition}[Restricted inverse step]\label{prop:restricted}
Let $0<q\leq1/100$, let $\|f\|_2=1$, and suppose
\begin{align}
  \E_S(f)&\geq1-q,\label{eq:restricted-primal}\\
  \E_T(\widehat{P_Sf})
    &\geq\left(1-\frac{q^2}{16}\right)N\,\E_S(f),
      \label{eq:restricted-fourier}\\
  |S|\,|T|&\leq\left(1+\frac{q^2}{16}\right)N.
      \label{eq:restricted-product}
\end{align}
Then there are $H,a,b$ satisfying
\eqref{eq:S-shape} and \eqref{eq:T-shape}. Moreover, if
$R_{H,a,b}$ denotes orthogonal projection onto the line spanned by
\[
  x\longmapsto |H|^{-1/2}(-1)^{b\cdot x}\1_{a+H}(x),
\]
then
\begin{equation}\label{eq:rank-one}
  \|f-R_{H,a,b}f\|_2^2\leq132q.
\end{equation}
\end{proposition}

Here is the mechanism behind Proposition~\ref{prop:restricted}. Restrict the
Walsh matrix to rows $T$ and columns $S$. Splitting $P_Sf$ into real and
imaginary parts shows that one normalized real component $v:S\to\mathbb R$
retains the ratio in \eqref{eq:restricted-fourier}. An exact Frobenius
identity then bounds the residual between the restricted sign matrix
$((-1)^{s\cdot t})_{t,s}$ and a rank-one matrix formed from $v$ and its
restricted Walsh transform:
\[
  \mathrm{Res}\leq \frac{q^2}{8}N.
\]
Rounding the two factors to signs costs at least one unit of squared residual
at every mismatching matrix entry. Hence the number of sign mismatches is at
most $q^2N/8$.

The four-point identity
\[
 (-1)^{s\cdot t}(-1)^{s_0\cdot t}
 (-1)^{s\cdot t_0}(-1)^{s_0\cdot t_0}
 =(-1)^{(s+s_0)\cdot(t+t_0)}
\]
turns those entry mismatches into almost orthogonality. Averaging over
anchors gives $s_0\in S$ and $t_0\in T$ such that, for
\[
  A=S+s_0,\qquad B=T+t_0,
\]
at most $q^2|A||B|$ pairs $(a,b)\in A\times B$ have $a\cdot b=1$, while
\[
  |A||B|\geq(1-q)N.
\]
The following finite additive lemma is the structural engine.

\begin{lemma}[Saturated almost orthogonality]\label{lem:saturated}
Suppose $A,B\subseteq V$, $A\neq\varnothing$, $0\in B$,
\[
  |A||B|\geq(1-q)N,
\]
and at most $q^2|A||B|$ pairs in $A\times B$ are nonorthogonal. If
$0<q\leq1/100$, then some subspace $L\leq V$ satisfies
\[
  |A\triangle L^\perp|\leq5q|A|,
  \qquad
  |B\triangle L|\leq12q|B|.
\]
\end{lemma}

For completeness, the lemma is obtained by taking
\[
  X=\{b\in B:\ |\{a\in A:a\cdot b=1\}|\leq q|A|\}.
\]
Finite Markov gives $|X|\geq(1-q)|B|$. The inequality
$\operatorname{bad}(x+x')\leq\operatorname{bad}(x)+
\operatorname{bad}(x')$, Walsh Parseval, and the saturation bound imply
$|X+X|<\tfrac32|X|$. The sharp small-doubling fact at the threshold $3/2$
therefore identifies $X+X$ with the point set of a subspace $L$.
Parseval and $|L||L^\perp|=N$ give
\[
  |L|\leq(1+10q)|B|,
  \qquad |L^\perp|\leq(1+3q)|A|.
\]
Combining these size bounds with $X\subseteq B\cap L$ and counting
nonorthogonal pairs outside $L^\perp$ yields the two asserted symmetric
differences.

Translating back gives the set conclusions in
Proposition~\ref{prop:restricted}. To recover the vector, let $P$ be the
coordinate projector onto $a+H$ and let $Q$ be the Fourier-coordinate
projector onto $b+H^\perp$. Flatness estimates from the same restricted
Walsh residual give
\[
  \|f-Pf\|_2^2\leq12q,\qquad
  \|f-Qf\|_2^2\leq54q.
\]
The projectors commute, and their common range is the one-dimensional line
in Proposition~\ref{prop:restricted}. Consequently,
\[
 \|f-PQf\|_2^2
 \leq2\|f-Pf\|_2^2+2\|f-Qf\|_2^2
 \leq132q,
\]
which proves \eqref{eq:rank-one}.

\subsection{From two-sided concentration to the restricted hypothesis}

\begin{proof}[Proof of Theorem~\ref{thm:main}]
It is enough to use the weaker bounds
\[
  \E_S(f)\geq1-\eta,\qquad
  \E_T(\widehat f)\geq(1-\eta)N,\qquad
  |S||T|\leq(1+\theta)N,
\]
where
\[
  \eta=\frac{q^4}{1024},\qquad
  \theta=\frac{q^2}{16}.
\]
Set
\[
  G=\E_S(f),\qquad
  A=\E_T(\widehat{P_Sf}),\qquad
  e=f-P_Sf.
\]
Then $\|e\|_2^2=1-G\leq\eta$, so Parseval gives
$\|\widehat e\|_2^2\leq\eta N$. For every $r>0$, the pointwise Young
inequality
\[
 |z+w|^2\leq(1+r)|z|^2+(1+r^{-1})|w|^2
\]
gives
\[
 \E_T(\widehat f)
 \leq(1+r)A+(1+r^{-1})\eta N.
\]
Choose $r=q^2/32=\theta/2$. The exact budget identity
\[
  (1+r)(1-\theta)+(1+r^{-1})\eta=1-\eta
\]
and the bound $G\leq1$ show that the right side of Young's inequality would
be at most $(1-\eta)N$ if $A$ were at most
$(1-\theta)NG$. Comparing with the lower bound on
$\E_T(\widehat f)$ therefore forces
\[
  A\geq(1-\theta)NG.
\]
Also $\eta\leq q$, so $G\geq1-q$. Thus all hypotheses of
Proposition~\ref{prop:restricted} hold, and it supplies $H,a,b$ and the
rank-one residual \eqref{eq:rank-one}.

Let $v$ be the normalized vector spanning that line and put
$d=\langle v,f\rangle$. Orthogonal projection gives
\[
  \|f-dv\|_2^2=1-|d|^2.
\]
Choose $c=d/|d|$ when $d\neq0$, and $c=1$ otherwise. Since $|d|\leq1$,
\[
  \|f-cv\|_2^2=2-2|d|
  \leq2(1-|d|^2)
  =2\|f-dv\|_2^2
  \leq264q.
\]
The vector $cv$ is exactly \eqref{eq:coset-state}, completing the proof.
\end{proof}

\section{Quantum interpretation}\label{sec:quantum}

Regard $f$ as the computational-basis amplitudes of an $n$-qubit pure state
\[
  |\psi\rangle=\sum_{x\in V}f(x)|x\rangle.
\]
Equation \eqref{eq:primal-concentration} says that computational-basis
measurement lands in $S$ with probability at least $1-\varepsilon$.
Because the amplitudes of the state after $H^{\otimes n}$ are
$\widehat f(y)/\sqrt N$, equation \eqref{eq:dual-concentration} says that
Hadamard-basis measurement lands in $T$ with the same probability.

The state \eqref{eq:coset-state} is an affine CSS stabilizer state. For
$h\in H$, translation $X^h$ has eigenvalue $(-1)^{b\cdot h}$; for
$k\in H^\perp$, the phase operator $Z^k$ has eigenvalue
$(-1)^{k\cdot a}$. Put less formally, ``CSS stabilizer state'' here means a
uniform superposition over solutions of linear equations in the bits, with
a linear pattern of plus and minus signs. The theorem therefore says that
near-minimal
concentration in two complementary bases certifies proximity to a hidden
CSS/coset state, not merely that the two concentration sets look affine.

Let $|\phi\rangle$ denote the state in \eqref{eq:coset-state}. For pure
states,
\[
  \disttr(\psi,\phi)
  =\sqrt{1-|\langle\psi|\phi\rangle|^2}
  \leq\||\psi\rangle-|\phi\rangle\|_2
  \leq\sqrt{264q}.
\]
Equivalently, the infidelity is at most $264q$. Every measurement outcome
distribution for the two states is therefore within total variation
distance $\sqrt{264q}$. These bounds are operationally meaningful but
numerically coarse; improving the state constant is a natural next target.

\section{Comparison with the continuous theorem}\label{sec:continuous}

It is useful to compare Theorem~\ref{thm:main} with the familiar uncertainty
principle for a particle moving on the real line. This comparison explains
what kind of result we have proved, but the two theorems are not special
cases of one another.

Let $\psi:\mathbb R\to\C$ be a normalized wavefunction. Its position
distribution is $|\psi(x)|^2$, and its momentum distribution is the squared
magnitude of its Fourier transform. Write $\sigma_x$ and $\sigma_\xi$ for
their standard deviations. With the normalization
\[
  \widehat\psi(\xi)
  =(2\pi)^{-1/2}\int_{\mathbb R}e^{-ix\xi}\psi(x)\,dx,
\]
the Heisenberg inequality is
\begin{equation}\label{eq:heisenberg}
  \sigma_x\sigma_\xi\geq\frac12.
\end{equation}
Equality holds exactly for Gaussian wave packets, after allowing
translation, modulation, dilation, and global phase
\cite{FollandSitaram1997}.

There is also a robust inverse statement: if the left side of
\eqref{eq:heisenberg} is close to $1/2$, then the wavefunction is close to a
Gaussian after applying those symmetries. The following elementary
one-dimensional version makes this precise.

\begin{proposition}[Continuous Gaussian stability]\label{prop:continuous}
Suppose $\|\psi\|_2=1$, its position and momentum variances are finite, and
\[
  \sigma_x\sigma_\xi\leq\frac{1+\kappa}{2}
\]
for $\kappa\geq0$. Then there are
$x_0,\xi_0,\theta\in\mathbb R$ and $s>0$ such that the normalized Gaussian
wave packet
\[
  g(x)=e^{i\theta}e^{i\xi_0x}s^{1/2}\pi^{-1/4}
       e^{-s^2(x-x_0)^2/2}
\]
satisfies
\[
  \|\psi-g\|_2^2\leq\kappa.
\]
\end{proposition}

\begin{proof}
Translate and modulate $\psi$ so that the two means are zero, then dilate it
so that the position and momentum variances are equal. These operations do
not change their product. For the resulting function $u$,
\[
  \langle u,(-d^2/dx^2+x^2)u\rangle
  =2\sigma_x\sigma_\xi\leq1+\kappa.
\]
The harmonic oscillator $-d^2/dx^2+x^2$ has a Gaussian ground state with
eigenvalue $1$, and every orthogonal state has energy at least $3$. This
spectral gap implies that the squared component of $u$ perpendicular to the
Gaussian is at most $\kappa/2$. If $\alpha$ is the overlap with the
normalized ground-state Gaussian, choosing the best global phase gives
\[
  \|u-e^{i\arg\alpha}g_0\|_2^2
  =2-2|\alpha|
  \leq2(1-|\alpha|^2)
  \leq\kappa.
\]
Undoing the translation, modulation, and dilation produces the required
Gaussian packet.
\end{proof}

The proposition is a standard spectral-gap argument; much sharper and
higher-order stability results are part of an active literature
\cite{DoLamLu2026}. The analogy with our finite theorem is:
\begin{center}
\begin{tabular}{p{0.30\linewidth}p{0.29\linewidth}p{0.29\linewidth}}
\hline
 & Continuous Heisenberg & Boolean support theorem\\
\hline
Space & Real line & $n$-bit strings\\
Measure of spread & Variance & Concentration-set size\\
Exact minimizers & Gaussian packets & Modulated coset states\\
Main rigidity tool & Harmonic-oscillator gap & Additive combinatorics\\
Robust conclusion & Close to a Gaussian & Close to a coset state\\
\hline
\end{tabular}
\end{center}

The shared principle is an \emph{inverse theorem}: near equality in an
uncertainty inequality forces closeness to the family of exact equality
cases. The important difference is the geometry. Variance and smoothness on
$\mathbb R$ lead to Gaussians and a linear spectral argument. Cardinality
and XOR-addition on $\F_2^n$ lead to subspace cosets and a combinatorial
argument. Consequently, the continuous theorem neither proves nor directly
implies Theorem~\ref{thm:main}.

\section{Relation to earlier work}\label{sec:literature}

\paragraph{Exact support uncertainty.}
Donoho and Stark established the finite support-product inequality in the
signal-recovery setting \cite{DonohoStark1989}. Meshulam and others developed
sharper uncertainty inequalities that depend on the subgroup structure of
the finite abelian group \cite{Meshulam2006}. Equality cases and related
generalizations are treated explicitly by Bonami--Ghobber and
Feng--Hollmann--Xiang
\cite{BonamiGhobber2013,FengHollmannXiang2019}; see also the survey of
Wigderson--Wigderson \cite{WigdersonWigderson2021}. These results supply the
exact baseline: a character restricted to a subgroup coset is an extremizer.
They do not by themselves give the leakage-tolerant quantitative conclusion
of Theorem~\ref{thm:main}.

\paragraph{Miller's robust coefficient inequality.}
Miller \cite[Theorem 3.2 and Corollary 3.1]{Miller2026} defines, for a
nonzero function on an arbitrary finite abelian group, a uniformity
coefficient $\nu(f)$ and a linearity coefficient $\eta(f)$. His corollary is
\[
  \frac{|\supp f|\,|\supp\widehat f|}{|G|}
  \,\nu(f)\eta(f)\geq1.
\]
Since both coefficients are at most one, a support product near $|G|$
forces them near one. Equality of $\nu$ means that $|f|$ is constant on its
support; equality of $\eta$ forces the support to be a subgroup coset. This
is a broad and elegant robust diagnostic, valid for all finite abelian
groups.

Miller's more general Theorem~3.2 bounds the Fourier overlap of two arbitrary
functions using support sizes together with $\nu$ and $\eta$. In the setting
of this note, one may apply it to the normalized restriction of the state to
$S$ and the normalized restriction of its Fourier transform to $T$. Small
leakage makes the resulting Fourier overlap close to one. Thus Miller's
theorem already gives a coefficient-level description of the
concentration-set near-extremal regime, even when the literal supports of the
original state and its transform are full.

Theorem~\ref{thm:main} is best viewed as a quantitative refinement of this
picture for the Boolean cube. It adds the following near-equality estimates:
\begin{itemize}
  \item There are a subspace $H$ and shifts $a,b$ such that $S$ differs from
  $a+H$ in at most $5q|S|$ points and $T$ differs from $b+H^\perp$ in at
  most $12q|T|$ points. Thus the two sets are close to a pair of
  complementary cosets built from the same subspace.
  \item The squared $\ell_2$ distance between the original state and the
  corresponding modulated coset state is at most $264q$, after allowing a
  global phase. This controls the complex amplitudes and relative signs, not
  only the supports or measurement probabilities.
  \item These geometric and state-vector estimates, including their
  constants, have a complete Lean formalization.
\end{itemize}
These conclusions trade generality for specificity. Miller works over
arbitrary finite abelian groups. The present theorem uses the special
exponent-two geometry of $\F_2^n$, including a four-point sign identity and
a sharp small-doubling step.
The proof given here is self-contained and does not invoke Miller's theorem.
Conceptually, however, it develops quantitative geometric and phase rigidity in
a Boolean regime diagnosed by his coefficient inequality. We do not claim a
new general uncertainty principle.

\paragraph{Continuous stability.}
The Gaussian conclusion for the Heisenberg variance inequality is the
closest conceptual analogue, as discussed in
Section~\ref{sec:continuous}. It belongs to a different branch of
uncertainty theory: its notion of spread, symmetry group, extremizers, and
proof mechanism all differ from the finite support problem. We include the
comparison because it clarifies the term ``robust inverse uncertainty,'' not
because either result subsumes the other.

\section{Lean formalization and audit surface}\label{sec:lean}

The formalization uses Lean 4.32.1 and Mathlib. It is organized so that an
auditor need not read the implementation before understanding the claim.
\begin{center}
\begin{tabular}{p{0.34\linewidth}p{0.58\linewidth}}
\hline
File & Role\\
\hline
\texttt{Specification.lean} &
Mathlib-only definitions of the cube, mass, Walsh transform, cosets,
hypotheses, and conclusion.\\
\texttt{Theorem.lean} &
The public theorem, its displayed mathematical docstring, and a one-line
proof delegation.\\
\texttt{Internal/Bridge.lean} &
Conversions between the public definitions and implementation definitions.\\
\texttt{Robust.lean} &
Analytic reduction, set and projector estimates, phase normalization, and
the final constants.\\
\texttt{Analytic.lean} &
Parseval identities, restricted-Walsh residuals, mismatch counting, and the
four-point argument.\\
\texttt{Combinatorial.lean} &
Finite Markov, small doubling, and saturated almost orthogonality.\\
\texttt{Basic.lean} &
Boolean Fourier and hidden-coset algebra.\\
\hline
\end{tabular}
\end{center}

\paragraph{Artifact availability.}
The complete Lean 4 formalization corresponding to this note is available in
the public, versioned repository release
\cite{GeorgiouLeanArtifact2026}. The tag \texttt{v0.1.0} identifies the
formal artifact checked in this paper.

The public declaration is:
\begin{verbatim}
theorem robust_inverse_uncertainty_phase
    (h : NearExtremizer epsilon delta q S T f) :
    exists H a b c, HiddenCosetApproximation q S T f H a b c
\end{verbatim}
The project contains no proof placeholders, project-local axioms, opaque
declarations, unsafe declarations, or uses of \texttt{native\_decide}. The
trusted base remains the Lean kernel, the imported Mathlib definitions and
theorems, and the ordinary compiler/runtime assumptions needed to check the
files \cite{deMouraUllrich2021,Mathlib2020}.

\section{Discussion}

Relative to Miller's coefficient-level result, the added content is a pair
of quantitative distance estimates in the Boolean setting. For near
equality, the theorem bounds how many points separate $S$ and $T$ from a
pair of complementary cosets, and it bounds how far the full complex state
is from the corresponding coset state. Formalization is useful here because
the argument combines Fourier normalization, additive combinatorics,
projector algebra, and many explicit constants.

There are also clear limitations. The result is local, the dependence on
the concentration loss is fourth-root, the constants are coarse, and the
argument is specialized to exponent two. A sharper proof might improve the
Young-inequality loss, avoid repeated factor-two estimates in the projector
step, or replace mismatch rounding with a weighted stability argument.
Extending the concentration-set theorem to general finite abelian groups
would require a structural lemma adapted to their subgroup and phase
geometry.

\paragraph{AI assistance disclosure.}
An AI system assisted with proof exploration, Lean implementation,
expository drafting, and literature organization. The repository records
the checkable formal artifact. Mathematical claims and decisions about
authorship, novelty, and publication remain the responsibility of the human
author.

\bibliographystyle{alpha}
\bibliography{references}

\end{document}
